\documentclass[10pt,a4paper]{amsart}
\usepackage{amsmath,amssymb,amsthm,graphicx}
% For commutative diagrams we use Xy-pic
\usepackage[all]{xy}

% We use 2cell for 2-commutative diagrams.
\xyoption{2cell}
% Theorem environments.
%
\theoremstyle{plain}
\newtheorem{theorem}[subsection]{Theorem}
\newtheorem{proposition}[subsection]{Proposition}
\newtheorem{lemma}[subsection]{Lemma}

\theoremstyle{definition}
\newtheorem{definition}[subsection]{Definition}
\newtheorem{example}[subsection]{Example}
\newtheorem{exercise}[subsection]{Exercise}
\newtheorem{situation}[subsection]{Situation}

\theoremstyle{remark}
\newtheorem{remark}[subsection]{Remark}
\newtheorem{remarks}[subsection]{Remarks}

\numberwithin{equation}{subsection}

% Macros
%
\def\lim{\mathop{\mathrm{lim}}\nolimits}
\def\colim{\mathop{\mathrm{colim}}\nolimits}
\def\Spec{\mathop{\mathrm{Spec}}}
\def\Hom{\mathop{\mathrm{Hom}}\nolimits}
\def\Ext{\mathop{\mathrm{Ext}}\nolimits}
\def\SheafHom{\mathop{\mathcal{H}\!\mathit{om}}\nolimits}
\def\SheafExt{\mathop{\mathcal{E}\!\mathit{xt}}\nolimits}
\def\Sch{\mathit{Sch}}
\def\Mor{\mathop{\mathrm{Mor}}\nolimits}
\def\Ob{\mathop{\mathrm{Ob}}\nolimits}
\def\Sh{\mathop{\mathit{Sh}}\nolimits}
\def\NL{\mathop{N\!L}\nolimits}
\def\CH{\mathop{\mathrm{CH}}\nolimits}
\def\proetale{{pro\text{-}\acute{e}tale}}
\def\etale{{\acute{e}tale}}
\def\QCoh{\mathit{QCoh}}
\def\Ker{\mathop{\mathrm{Ker}}}
\def\Im{\mathop{\mathrm{Im}}}
\def\Coker{\mathop{\mathrm{Coker}}}
\def\Coim{\mathop{\mathrm{Coim}}}

% Boxtimes
%
\DeclareMathSymbol{\boxtimes}{\mathbin}{AMSa}{"02}

%
% Macros for moduli stacks/spaces
%
\def\QCohstack{\mathcal{QC}\!\mathit{oh}}
\def\Cohstack{\mathcal{C}\!\mathit{oh}}
\def\Spacesstack{\mathcal{S}\!\mathit{paces}}
\def\Quotfunctor{\mathrm{Quot}}
\def\Hilbfunctor{\mathrm{Hilb}}
\def\Curvesstack{\mathcal{C}\!\mathit{urves}}
\def\Polarizedstack{\mathcal{P}\!\mathit{olarized}}
\def\Complexesstack{\mathcal{C}\!\mathit{omplexes}}
% \Pic is the operator that assigns to X its picard group, usage \Pic(X)
% \Picardstack_{X/B} denotes the Picard stack of X over B
% \Picardfunctor_{X/B} denotes the Picard functor of X over B
\def\Pic{\mathop{\mathrm{Pic}}\nolimits}
\def\Picardstack{\mathcal{P}\!\mathit{ic}}
\def\Picardfunctor{\mathrm{Pic}}
\def\Deformationcategory{\mathcal{D}\!\mathit{ef}}
\usepackage{cleveref}
\crefname{originaltheorem}{original theorem}{original theorems}
\Crefname{originaltheorem}{Original theorem}{Original theorems}
\makeatletter
\newlabel{theorem-henselization}{{03QL}{source}{}{sourcepointer-1}{}}
\newlabel{theorem-henselization@cref}{{[originaltheorem][0][]03QL}{[1][0][]source}{}{}{}}
\makeatother
\crefname{originallemma}{original lemma}{original lemmas}
\Crefname{originallemma}{Original lemma}{Original lemmas}
\makeatletter
\newlabel{algebra-lemma-strict-henselization-different}{{04GW}{source}{}{sourcepointer-2}{}}
\newlabel{algebra-lemma-strict-henselization-different@cref}{{[originallemma][0][]04GW}{[1][0][]source}{}{}{}}
\makeatother
\title{The stalk of the structure sheaf in the etale topology is the strict henselization}
\author{The Stacks Project Authors}
\date{stacks-project:etale-cohomology.tex:local-snapshot-20261011; GNU FDL 1.2 or later}
\begin{document}
\maketitle
\noindent\textit{Editorial note.} All delivered mathematical text is the author's own, copied from the pinned Stacks Project chapter file etale-cohomology.tex; no statement or proof step was rewritten, reordered or removed.\par\medskip
\noindent\textit{Editorial note.} The proof is the author's complete argument; the two colimit descriptions are written out and identified.\par\medskip
\noindent\textit{Editorial note.} Every labelled reference occurring inside the counted spans is carried in assembly.external\_references as a non-load-bearing navigation pointer (same-chapter ones at their definition site, cross-chapter ones at the invocation inside the counted span), so no load-bearing mathematics is replaced by a stub. No bibliography citation occurs inside the counted spans.\par\medskip
\noindent\textit{Editorial note.} The checked chapter list for same-name collisions is the set of stacks chapters staged in this repository: chow, cohomology, divisors, etale-cohomology, flat, more-algebra, more-morphisms, morphisms, spaces-more-morphisms, varieties together with this card's chapter etale-cohomology.\par\medskip
\noindent\textit{Editorial note.} The preamble is the chapter's own native preamble excerpt already used for this source: the Xy-pic setup (native preamble lines 14-18) and the native theorem declarations with the mathematical macros (native preamble lines 154-217).\par\medskip
\noindent\textit{Allowed standard prerequisites (editorial).}\begin{itemize}
\item Strict henselization of a local ring and its description by a colimit of etale algebras. Same chapter, the theorem describing the strict henselization of a local ring as a filtered colimit of etale algebras. Supplies the left hand side of the asserted identification.
\item Stalk of the structure sheaf at a geometric point. Same chapter, the description of the stalk of O\_S at a geometric point as a filtered colimit over etale neighbourhoods. Supplies the right hand side and the comparison with the colimit description of the strict henselization.
\end{itemize}\medskip
\ifdefined\hypertarget\hypertarget{sourcepointer-1}{}\fi
\noindent\textit{Source reference.} Original theorem 03QL: sources/stacks-etale-cohomology/pinned-original-source.tex, source lines 4216--4216. Navigation pointer invoked at proof line 4237; the counted span contains the pointer only and the cited result is not reproduced inside it, so the pointer is registered as a non-load-bearing navigation stub in the established form for Stacks material. No bibliography citation is involved.\par\medskip
\ifdefined\hypertarget\hypertarget{sourcepointer-2}{}\fi
\noindent\textit{Source reference.} Original lemma 04GW: sources/stacks-etale-cohomology/pinned-original-source.tex, source lines 4308--4308. Navigation pointer invoked at proof line 4308; the counted span contains the pointer only and the cited result is not reproduced inside it, so the pointer is registered as a non-load-bearing navigation stub in the established form for Stacks material. No bibliography citation is involved.\par\medskip
\begin{lemma}
\label{lemma-describe-etale-local-ring}

Let $S$ be a scheme.
Let $\overline{s}$ be a geometric point of $S$ lying over $s \in S$.
Let $\kappa = \kappa(s)$ and let
$\kappa \subset \kappa^{sep} \subset \kappa(\overline{s})$ denote
the separable algebraic closure of $\kappa$ in $\kappa(\overline{s})$.
Then there is a canonical identification
$$
(\mathcal{O}_{S, s})^{sh}
\cong
(\mathcal{O}_S)_{\overline{s}}
$$
where the left hand side is the strict henselization of the local ring
$\mathcal{O}_{S, s}$ as described in
Theorem \ref{theorem-henselization}
and right hand side is the stalk of the structure sheaf
$\mathcal{O}_S$ on $S_\etale$ at
the geometric point $\overline{s}$.
\end{lemma}

\begin{proof}
Let $\Spec(A) \subset S$ be an affine neighbourhood of $s$.
Let $\mathfrak p \subset A$ be the prime ideal corresponding to $s$.
With these choices we have canonical isomorphisms
$\mathcal{O}_{S, s} = A_{\mathfrak p}$ and $\kappa(s) = \kappa(\mathfrak p)$.
Thus we have
$\kappa(\mathfrak p) \subset \kappa^{sep} \subset \kappa(\overline{s})$.
Recall that
$$
(\mathcal{O}_S)_{\overline{s}} =
\colim_{(U, \overline{u})} \mathcal{O}(U)
$$
where the limit is over the \'etale neighbourhoods of $(S, \overline{s})$.
A cofinal system is given by those \'etale neighbourhoods $(U, \overline{u})$
such that $U$ is affine and $U \to S$ factors through $\Spec(A)$.
In other words, we see that
$$
(\mathcal{O}_S)_{\overline{s}} = \colim_{(B, \mathfrak q, \phi)} B
$$
where the colimit is over \'etale $A$-algebras $B$ endowed with a prime
$\mathfrak q$ lying over $\mathfrak p$ and a
$\kappa(\mathfrak p)$-algebra map
$\phi : \kappa(\mathfrak q) \to \kappa(\overline{s})$.
Note that since $\kappa(\mathfrak q)$ is finite separable over
$\kappa(\mathfrak p)$ the image of $\phi$ is contained in $\kappa^{sep}$.
Via these translations the result of the lemma is equivalent
to the result of
Algebra, Lemma \ref{algebra-lemma-strict-henselization-different}.
\end{proof}

\end{document}
